\section{Machine Design Concepts, Core Materials, Copper and Core Losses}

\subsection{Magnetic Core Materials and Laminations (Solid Mild Steel vs Laminated Silicon Steel, $k_{\mathrm{fe}}$)}

\subsubsection{Solid Mild Steel vs Laminated Silicon Steel: DC vs AC Flux and Eddy Currents}
The choice of magnetic core material is governed by whether the magnetic flux is static (DC) or time-varying (AC):
\begin{itemize}
  \item \textbf{Solid mild steel (cast steel):} Used strictly for \textbf{unidirectional (DC) magnetic flux}, such as the rotor bodies, rotor spider, and field poles of synchronous machines excited by direct current $I_f$. Solid steel provides high mechanical tensile strength and high saturation flux density ($B_{\mathrm{sat}} \approx 1.8\text{--}2.0\,\mathrm{T}$) at low manufacturing cost.
  \item \textbf{Fatal trap with solid steel under AC flux:} If solid steel is subjected to alternating flux, macroscopic eddy currents circulate freely across the large cross-section of low-resistivity steel, generating excessive $I^2 R$ loss, severe core overheating, and thermal breakdown.
  \item \textbf{Laminated silicon steel (electrical steel sheets):} Used for all \textbf{alternating (AC) magnetic flux paths}, including stator cores of all AC machines, induction motor rotor cores, and pole shoes of synchronous machines facing slotted stators.
  \item \textbf{Role of silicon alloy (2\%--4\%):} Adding silicon increases the electrical resistivity $\rho_{\mathrm{fe}}$ of the steel by a factor of 4 to 5, which drastically suppresses eddy-current flow while preserving high magnetic permeability.
\end{itemize}

\subsubsection{Lamination Stacking Factor ($k_{\mathrm{fe}}$) and Net Core Iron Area ($A_{\mathrm{core}} = k_{\mathrm{fe}} A_{\mathrm{gross}}$)}
Laminated cores are constructed from thin insulated stampings (typically $d_{\mathrm{lam}} = 0.35\text{--}0.50\,\mathrm{mm}$ thick).
Due to insulating varnish, surface oxide, and imperfect compression, the active iron volume is slightly less than the gross core volume.
The stacking / packing factor $k_{\mathrm{fe}}$ is defined as:
\[
k_{\mathrm{fe}} = \frac{A_{\mathrm{core}}}{A_{\mathrm{gross}}} = \frac{V_{\mathrm{iron}}}{V_{\mathrm{gross}}} \approx 0.90\text{--}0.95.
\]
The net magnetic iron cross-sectional area $A_{\mathrm{core}}$ available for carrying magnetic flux is:
\[
A_{\mathrm{core}} = k_{\mathrm{fe}} A_{\mathrm{gross}} = k_{\mathrm{fe}} L_c b_{\mathrm{core}}.
\]

\subsection{Core Losses: Eddy-Current and Hysteresis Loss Mechanisms ($P_{\mathrm{core}}$, $P_{\mathrm{eddy}}$, $P_{\mathrm{hyst}}$)}

\subsubsection{Eddy-Current Loss Dependence ($P_{\mathrm{eddy}} \propto f^2 B_{\mathrm{core}}^2 d_{\mathrm{lam}}^2 / \rho_{\mathrm{fe}}$)}
Alternating magnetic flux induces circulating currents (eddy currents) in the conducting core laminations.
The volumetric eddy-current power loss is governed by:
\[
P_{\mathrm{eddy}} = \frac{\pi^2}{6 \rho_{\mathrm{fe}}} f^2 B_{\mathrm{core}}^2 d_{\mathrm{lam}}^2 \cdot V_{\mathrm{core}} = k_e\,f^2\,B_{\mathrm{core}}^2\,d_{\mathrm{lam}}^2\quad[\mathrm{W}].
\]
\textbf{Key design dependencies:}
\begin{itemize}
  \item Quadratic in frequency: $P_{\mathrm{eddy}} \propto f^2$.
  \item Quadratic in core flux density: $P_{\mathrm{eddy}} \propto B_{\mathrm{core}}^2$.
  \item Quadratic in lamination thickness: $P_{\mathrm{eddy}} \propto d_{\mathrm{lam}}^2$. Halving sheet thickness from $0.50\,\mathrm{mm}$ to $0.25\,\mathrm{mm}$ cuts eddy-current loss to $1/4$ ($25\%$).
  \item Inversely proportional to steel resistivity $\rho_{\mathrm{fe}}$: Adding silicon increases $\rho_{\mathrm{fe}}$, reducing $P_{\mathrm{eddy}}$.
\end{itemize}

\subsubsection{Hysteresis Loss and Steinmetz Empirical Relation ($P_{\mathrm{hyst}} \propto f B_{\mathrm{core}}^n$)}
Hysteresis loss represents the energy dissipated per cycle in traversing the ferromagnetic $B$--$H$ hysteresis loop.
According to the Steinmetz empirical formula:
\[
P_{\mathrm{hyst}} = k_h\,f\,B_{\mathrm{core}}^n\cdot V_{\mathrm{core}}\quad[\mathrm{W}],
\]
where $n \approx 1.6\text{--}2.0$ (Steinmetz exponent; commonly approximated as $n=2$ in linear scaling exercises).
\textbf{Key design dependencies:}
\begin{itemize}
  \item Linear in frequency: $P_{\mathrm{hyst}} \propto f$.
  \item Power law in peak flux density: $P_{\mathrm{hyst}} \propto B_{\mathrm{core}}^2$.
\end{itemize}

\subsubsection{Total Core Loss Proportionality to Voltage ($P_{\mathrm{core}} \propto B_{\mathrm{core}}^2 \propto V_{\mathrm{ph}}^2$)}
Total iron core loss is the sum of eddy-current and hysteresis components:
\[
P_{\mathrm{core}} = P_{\mathrm{eddy}} + P_{\mathrm{hyst}}.
\]
From Faraday's law, induced phase EMF is $E_{\mathrm{ph}} \approx V_{\mathrm{ph}} \approx 4.44\,K_w T_{\mathrm{ph}} \Phi_m f \propto f B_{\mathrm{core}}$.
At fixed electrical supply frequency $f$:
\[
B_{\mathrm{core}} \propto \Phi_m \propto V_{\mathrm{ph}}
\implies
P_{\mathrm{core}} \propto B_{\mathrm{core}}^2 \propto V_{\mathrm{ph}}^2.
\]

\subsection{Stator Copper Loss and Thermal Trade-Offs ($P_{\mathrm{Cu}}$, $I_{\mathrm{ph}}$, $R_{\mathrm{ph}}$)}

\subsubsection{Three-Phase Copper Loss Formulation ($P_{\mathrm{Cu}} = 3 I_{\mathrm{ph}}^2 R_{\mathrm{ph}} \propto I_{\mathrm{ph}}^2$)}
Total stator winding Joule heating (copper loss) across all three phases is:
\[
P_{\mathrm{Cu}} = 3 I_{\mathrm{ph}}^2 R_{\mathrm{ph}}\quad[\mathrm{W}].
\]
Copper loss depends quadratically on operating phase current $I_{\mathrm{ph}}^2$, making it strongly load-dependent.
For synchronous machines, the DC field winding copper loss is:
\[
P_{\mathrm{Cu,f}} = I_f^2 R_f\quad[\mathrm{W}].
\]

\subsubsection{Current Density Sizing and Thermal Balance}
Conductor current density $J_{\mathrm{cond}} = I_c / A_{\mathrm{cond}}$ governs conductor cross-section and winding temperature:
\begin{itemize}
  \item If $J_{\mathrm{cond}}$ is chosen too high: Slot area $A_{\mathrm{slot}}$ is small (compact core), but phase resistance $R_{\mathrm{ph}}$ and copper loss $P_{\mathrm{Cu}}$ rise rapidly, causing excessive thermal rise and thermal insulation aging.
  \item If $J_{\mathrm{cond}}$ is chosen too low: Conductor area $A_{\mathrm{cond}}$ becomes large, requiring deeper or wider slots. This forces teeth to be thinner (causing iron tooth saturation) or increases overall machine frame size and copper raw material cost.
\end{itemize}

\subsection{Loss Balancing, Sizing Trade-Offs and Efficiency Optimization ($P_{\mathrm{Cu}} \approx P_{\mathrm{core}}$, $\eta_{\max}$)}

\subsubsection{Loss Balancing Principle for Maximum Full-Load Efficiency}
Machine overall efficiency $\eta$ is:
\[
\eta = \frac{P_{\mathrm{out}}}{P_{\mathrm{in}}} = \frac{P_{\mathrm{out}}}{P_{\mathrm{out}} + P_{\mathrm{Cu}} + P_{\mathrm{core}} + P_{\mathrm{mech,loss}}}.
\]
Since copper loss $P_{\mathrm{Cu}} \propto I_{\mathrm{ph}}^2$ is variable and core loss $P_{\mathrm{core}} \propto V_{\mathrm{ph}}^2$ is constant with respect to load current, mathematical differentiation reveals that \textbf{efficiency peaks when variable copper loss equals constant core loss}:
\[
P_{\mathrm{Cu}} \approx P_{\mathrm{core}}\quad\text{at rated design operating point}.
\]

\subsubsection{Economic and Volume Trade-Offs: Excess Copper vs Excess Iron}
\begin{itemize}
  \item \textbf{Excess copper (over-sized conductors):} Attempting to drive copper loss close to zero requires massive slot areas. Stator teeth become exceedingly narrow and suffer intense magnetic saturation, or the outer stator diameter must be enlarged, severely increasing machine size, weight, and copper material expenditure.
  \item \textbf{Excess iron (over-sized magnetic core):} Attempting to minimize core loss requires operating at very low magnetic flux density $B_{\mathrm{core}}$. This requires massive core cross-sections $A_{\mathrm{core}}$ and long stack length $L_c$, leading to an excessively heavy, costly machine frame.
  \item \textbf{Balanced design:} Sizing the machine so that $P_{\mathrm{Cu}} \approx P_{\mathrm{core}}$ at rated load achieves the optimum balance between physical size, material cost, and electrical efficiency.
\end{itemize}

\subsubsection{Partial-Load Efficiency and Voltage Reduction Strategy ($V_{\mathrm{ph}} \downarrow$ as $P_{\mathrm{mech}} \downarrow$)}
When an electric motor operates below rated shaft load ($P_{\mathrm{mech}} < P_{\mathrm{rated}}$):
\begin{enumerate}
  \item Motor current $I_{\mathrm{ph}}$ drops, so copper loss falls rapidly ($P_{\mathrm{Cu}} \propto I_{\mathrm{ph}}^2$).
  \item If terminal voltage $V_{\mathrm{ph}}$ is held constant at rated voltage, core loss $P_{\mathrm{core}} \propto V_{\mathrm{ph}}^2$ remains completely unchanged.
  \item As a result, at partial load $P_{\mathrm{core}} \gg P_{\mathrm{Cu}}$, causing overall motor efficiency $\eta$ to decline sharply.
  \item \textbf{Efficiency optimization via voltage scaling:} By reducing terminal voltage $V_{\mathrm{ph}}$ under partial load (e.g.\ via an inverter or soft-starter), iron core loss is reduced ($P_{\mathrm{core}} \propto V_{\mathrm{ph}}^2$) until $P_{\mathrm{core}} \approx P_{\mathrm{Cu}}$, restoring high operating efficiency.
\end{enumerate}

\subsection{Magnetic Core Saturation: Magnetizing Surge and Power Factor Degradation ($B_{\mathrm{core}}$, $I_\mu$, $\cos\phi$)}

\subsubsection{Non-Linear $B$--$H$ Characteristics and Core Saturation Limits}
Real magnetic sheet steels have a non-linear magnetization curve:
\begin{itemize}
  \item \textbf{Linear region ($B < 1.4\,\mathrm{T}$):} Relative permeability is high ($\mu_r \approx 2000\text{--}8000$). Magnetic reluctance of iron is nearly negligible compared to the air gap.
  \item \textbf{Knee / saturation region ($B > 1.6\text{--}1.8\,\mathrm{T}$):} Iron saturates and differential permeability drops drastically towards $\mu_0$. The iron segments begin to demand enormous MMF.
\end{itemize}

\subsubsection{Consequences of Over-Fluxing: No-Load Current $I_0 \uparrow$, Harmonics, and Power Factor $\cos\phi \downarrow$}
When core flux density $B_{\mathrm{core}}$ exceeds design limits:
\begin{enumerate}
  \item \textbf{Sharp surge in magnetizing current $I_\mu$:} The required ampere-turns $\oint \mathbf{H}\cdot\mathrm{d}\mathbf{l}$ jump steeply, forcing magnetizing current $I_\mu$ to rise by $100\%\text{--}400\%$.
  \item \textbf{No-load current increases:} Total no-load current $I_0 = \sqrt{I_w^2 + I_\mu^2}$ rises correspondingly, primarily due to reactive component $I_\mu$.
  \item \textbf{Severe degradation of power factor $\cos\phi$:} Because magnetizing current is purely reactive (lags applied voltage by $90^\circ$), a large $I_\mu$ drives the overall power factor $\cos\phi$ down sharply.
  \item \textbf{Elevated core losses and overheating:} $P_{\mathrm{core}} \propto B_{\mathrm{core}}^2$ increases, causing thermal hotspots in teeth and stator yoke.
  \item \textbf{Harmonic generation:} Non-linear $B$--$H$ saturation distorts sinusoidal flux waves, generating pronounced 3rd and other odd time-harmonic currents or voltages.
\end{enumerate}

\subsection{Exam Recipe: Evaluating the Consequences of Magnetic Core Saturation}
\textbf{Use when:} Problem asks what happens if machine flux density $B$, core loading, or applied voltage $V$ is raised beyond rated levels into iron saturation.

\textbf{Steps:}
\begin{enumerate}
  \item \textbf{Identify flux density increase:} Note that $B_{\mathrm{core}} \propto V_{\mathrm{ph}}/f$. Higher voltage or lower frequency causes over-fluxing.
  \item \textbf{State core loss increase:} Core losses increase quadratically: $P_{\mathrm{core}} = P_{\mathrm{eddy}} + P_{\mathrm{hyst}} \propto B_{\mathrm{core}}^2$.
  \item \textbf{State magnetizing current surge:} In the saturated region of the $B$--$H$ curve, $\mu_r$ plummets; hence required magnetizing MMF explodes, causing a dramatic surge in magnetizing current $I_\mu$ and no-load current $I_0$.
  \item \textbf{State power factor penalty:} Because the surge in $I_0$ is almost entirely reactive current $I_\mu$, the operating power factor $\cos\phi$ deteriorates significantly.
  \item \textbf{State waveform distortion:} Saturation introduces harmonic flux and current components (especially 3rd harmonic).
\end{enumerate}

\textbf{Fast checks:} Saturation never produces ``free torque'' or higher efficiency; it always elevates losses, increases reactive current, and worsens power factor $\cos\phi$.

\subsection{Exam Recipe: Partial-Load Voltage Reduction for Maintaining High Efficiency}
\textbf{Use when:} Problem asks how to improve or optimize motor efficiency when operating continuously at reduced mechanical load ($P_{\mathrm{mech}} < P_{\mathrm{rated}}$).

\textbf{Steps:}
\begin{enumerate}
  \item \textbf{Identify the loss imbalance:} At partial load, $I_{\mathrm{ph}}$ drops, so $P_{\mathrm{Cu}} = 3 I_{\mathrm{ph}}^2 R_{\mathrm{ph}}$ drops significantly, while constant voltage keeps $P_{\mathrm{core}} \propto V_{\mathrm{ph}}^2$ unnecessarily high ($P_{\mathrm{core}} \gg P_{\mathrm{Cu}}$).
  \item \textbf{Apply the efficiency criterion:} Maximum efficiency requires $P_{\mathrm{Cu}} \approx P_{\mathrm{core}}$.
  \item \textbf{Reduce terminal voltage $V_{\mathrm{ph}}$:} Lowering terminal voltage proportionally reduces core flux density $B_{\mathrm{core}} \propto V_{\mathrm{ph}}$ and therefore core loss $P_{\mathrm{core}} \propto V_{\mathrm{ph}}^2$.
  \item \textbf{Evaluate the new balance:} Re-balancing $P_{\mathrm{core}} \approx P_{\mathrm{Cu}}$ at the lower operating point restores peak operating efficiency $\eta$.
\end{enumerate}

\subsection{Common Traps and Quick Sanity Checks in Machine Design and Losses}
\begin{itemize}
  \item \textbf{Solid steel vs laminated core:} Solid mild steel is suitable ONLY for DC flux paths (synchronous machine rotor bodies). Never use solid steel for AC stators or induction rotors due to catastrophic eddy-current overheating.
  \item \textbf{Lamination thickness dependence:} Eddy current loss is proportional to $d_{\mathrm{lam}}^2$, NOT $d_{\mathrm{lam}}$. Changing thickness from $0.5\,\mathrm{mm}$ to $0.35\,\mathrm{mm}$ reduces $P_{\mathrm{eddy}}$ by a factor of $(0.35/0.50)^2 \approx 0.49$ (a $51\%$ reduction).
  \item \textbf{Core loss vs load current:} At constant supply voltage and frequency, core loss $P_{\mathrm{core}}$ is virtually independent of load current $I_{\mathrm{ph}}$. Copper loss $P_{\mathrm{Cu}}$ varies with $I_{\mathrm{ph}}^2$.
  \item \textbf{No-load current composition:} No-load current $I_0$ consists of core-loss active current $I_w$ and magnetizing reactive current $I_\mu$ ($I_0 = \sqrt{I_w^2 + I_\mu^2}$). In AC machines, $I_\mu \gg I_w$, so $I_0$ is predominantly reactive.
  \item \textbf{Over-saturation fallacy:} Higher $B$ does not automatically increase motor output safely; past the saturation knee, iron saturation leads to heavy overheating and severely lagging power factor.
\end{itemize}
